% Part: computability
% Chapter: recursive-functions
% Section: notations-pr-functions

\documentclass[../../../include/open-logic-section]{subfiles}

\begin{document}

\olfileid[fr]{cmp}{rec}{not}
\olsection{Notations de la récurrence primitive}

La définition inductive précise des fonctions récursives primitives
permet de les décrire systématiquement. On peut, par exemple,
attribuer une « notation » à chacune d'elles. Utilisons les
symboles $\Zero$, $\Succ$ et $\Proj{n}{i}$ pour désigner
respectivement la fonction nulle, la fonction successeur et les
projections. Supposons maintenant que $h$ soit définie par
composition à partir d'une fonction~$f$ à $k$ arguments et des
fonctions $g_0,\dots,g_{k-1}$ à $n$ arguments. Si les notations
attribuées respectivement à $f,g_0,\dots,g_{k-1}$ sont
$F,G_0,\dots,G_{k-1}$, nous pouvons, à l'aide d'un nouveau
symbole~$\fn{Comp}_{k,n}$, noter $h$ par
$\fn{Comp}_{k,n}[F,G_0,\dots,G_{k-1}]$.

On emploie des notations analogues pour les fonctions définies
par récurrence primitive. Supposons que la fonction~$h$ à
$k+1$ arguments soit définie par récurrence primitive à partir
de la fonction~$f$ à $k$ arguments et de la fonction~$g$ à
$k+2$ arguments, auxquelles sont respectivement attribuées les
notations $F$ et~$G$. La notation de~$h$ est alors
$\fn{Rec}_k[F,G]$.

Rappelons que l'addition est définie par récurrence primitive
comme suit :
\begin{align*}
  \Add(x_0, 0) & = \Proj{1}{0}(x_0) = x_0\\
  \Add(x_0, y+1) & = \Succ(\Proj{3}{2}(x_0, y, \Add(x_0, y))) = \Add(x_0, y) +1
\end{align*}
Le rôle de~$f$ est tenu par $\Proj{1}{0}$, et celui de~$g$
par $\Succ(\Proj{3}{2}(x_0,y,z))$. Cette dernière fonction
porte la notation $\fn{Comp}_{1,3}[\Succ,\Proj{3}{2}]$ :
elle résulte de la composition de la fonction~$\Succ$ à un
argument avec la fonction~$\Proj{3}{2}$ à trois arguments.
L'addition reçoit donc la notation
\[
\fn{Rec}_1[\Proj{1}{0},\fn{Comp}_{1,3}[\Succ,\Proj{3}{2}]].
\]
Ces notations sont parfois utiles, notamment pour énumérer les
fonctions récursives primitives.

\begin{prob}
  Donner la notation récursive primitive complète de~$\Mult$.
\end{prob}

\end{document}
